\documentclass[10pt,a4paper]{amsart}
\usepackage[margin=19mm]{geometry}
\usepackage{amsmath,amssymb,mathtools}
\usepackage[scr=rsfs,cal=euler]{mathalfa}
\usepackage{xfrac}
\usepackage[unicode,hidelinks,bookmarks=false]{hyperref}
\newcommand{\ie}{i.e.\ }
\newcommand{\cf}{cf.\ }
\newcommand{\resp}{respectively\ }
\newcommand{\Subsection}{\subsection}
\providecommand{\nobreakdash}{\nobreak\hbox{-}}
\setlength{\emergencystretch}{2em}
\multlinegap0pt
\usepackage{amssymb}
\usepackage{mathtools}
\usepackage{bm}
\usepackage[shortlabels]{enumitem}
\newtheorem{lemma}{Lemma}
\newtheorem{proposition}{Proposition}
\newcommand{\Hess}{\mathcal H}
\newcommand{\LOC}{\operatorname{LOC}}
\newcommand{\Sd}{\mathbb S^d}
\newcommand{\Spd}{\mathbb S^d_+}
\newcommand{\dist}{\operatorname{dist}}
\title{Local Nonconvexity Indices for \(C^{1,1}\) Functions via Generalized Hessians}
\author{Marina Palaisti}
\date{Source: arXiv:2605.08491v1 (CC BY 4.0)}
\begin{document}
\maketitle

\noindent\emph{Source and licence.} Local Nonconvexity Indices for \(C^{1,1}\) Functions via Generalized Hessians. Version arXiv:2605.08491v1 (CC BY 4.0), distributed by arXiv under the Creative Commons Attribution 4.0 International licence, http://creativecommons.org/licenses/by/4.0/, as recorded in the arXiv licence metadata for this exact version and shown on the abstract page https://arxiv.org/abs/2605.08491v1. Full licence notice: \texttt{licenses/arxiv-2605.08491-CC-BY-4.0.txt}.\par\smallskip

\noindent\textit{Editorial scope.} Counted unit: the article's proposition prop:sum (source0.tex lines 467--477). The article's complete original proof of it is delivered with it (source0.tex lines 479--524, one \texttt{proof} environment of 46 lines). No mathematics is written, repaired or re-derived: the statement and the proof are the article's own lines, in the article's own notation, and no retained line is substituted or re-flowed.  Retained uncounted as the source-local context the proof needs: lines 147--159.  Nothing was omitted from any retained span, so the package carries no omission record.  No retained line contains a citation, so the wrapper carries no bibliography and no bibliography entry.  No retained line contains a figure environment, an image-inclusion command or drawing code, so no image is delivered and the package declares no asset dependency.  Editorial formatting only, in addition: an AMS \texttt{amsart} preamble that restates the article's own declarations and macros used by the retained lines, namely the environments \texttt{lemma}, \texttt{proposition} on separate counters, together with the standard packages those lines need.  The retained environments print in this document as Lemma 1, Proposition 1 in order of appearance, whereas the article prints the same items with the numbering of its own full document.  The complete native source of the article as distributed in the arXiv e-print for this exact version is bundled unchanged as \texttt{sources/200157/source0.tex}.
	\begin{lemma}[Distance from the positive semidefinite cone]
		\label{lem:distance}
		For every \(Q\in\Sd\),
		\[
		\dist_*(Q,\Spd)
		:=
		\inf_{M\in\Spd}\|Q-M\|_*
		=
		\|Q^-\|_*
		=
		\ell(Q).
		\]
	\end{lemma}

	\begin{proposition}[Upper endpoint subadditivity]
		\label{prop:sum}
		Let \(h_1,h_2\in C^{1,1}_{\mathrm{loc}}(G)\) and set \(h=h_1+h_2\). Then, for every \(x\in G\),
		\[
		\overline{\LOC}(h;x)
		\leq
		\overline{\LOC}(h_1;x)
		+
		\overline{\LOC}(h_2;x).
		\]
	\end{proposition}

	\begin{proof}
		Since
		\[
		\nabla h=\nabla h_1+\nabla h_2,
		\]
		the standard Clarke generalized Jacobian sum rule for locally Lipschitz mappings gives
		\[
		\Hess(h;x)
		\subset
		\Hess(h_1;x)+\Hess(h_2;x),
		\]
		where the sum on the right is the Minkowski sum of subsets of \(\Sd\).
		
		Let \(Q\in\Hess(h;x)\). Then \(Q=Q_1+Q_2\) for some
		\[
		Q_i\in\Hess(h_i;x),
		\qquad i=1,2.
		\]
		Using Lemma~\ref{lem:distance},
		\[
		\ell(Q)
		=
		\dist_*(Q,\Spd).
		\]
		For arbitrary \(M_1,M_2\in\Spd\), we have \(M_1+M_2\in\Spd\), and therefore
		\[
		\ell(Q_1+Q_2)
		\leq
		\|Q_1+Q_2-(M_1+M_2)\|_*
		\leq
		\|Q_1-M_1\|_*+\|Q_2-M_2\|_*.
		\]
		Taking infima over \(M_1,M_2\in\Spd\) yields
		\[
		\ell(Q_1+Q_2)
		\leq
		\ell(Q_1)+\ell(Q_2).
		\]
		Thus
		\[
		\ell(Q)
		\leq
		\overline{\LOC}(h_1;x)+\overline{\LOC}(h_2;x).
		\]
		Taking the supremum over \(Q\in\Hess(h;x)\) proves the result.
	\end{proof}

\end{document}
